\Def Медианой массива называют его $\frac{N}{2}$-ю порядковую статистику. Получается, оптимальнее всего брать медиану в качестве pivot.

Нам требуется описать алгоритм поиска медианы за линейное время, но мы будем решать более общую задачу поиска k-й порядковой статистики за линейное время, которая покрывает задачу поиска медианы

\Alg Поиск медианы медиан
\begin{enumerate}
\item Разбить массив на блоки по пять.
\item Отсортировать блоки (можно пузырьком).
\item Выбрать медианы блоков.
\item Для массива медиан найдем его медиану. Найденный элемент назовем медианой медиан.
\end{enumerate}

\subsection*{Детерминированный алгоритм поиска k-й порядковой статистики}

\Statement Медиана медиан гарантированно делит массив в соотношении не хуже чем 3 : 7.

\Proof Cначала определим нижнюю границу для количества элементов, превышающих по величине опорный элемент x. В общем случае как минимум половина медиан, найденных на втором шаге, больше или равны медианы медиан x.
\\Таким образом, как минимум $\frac{N}{10}$ групп содержат по 3 элемента, превышающих величину x, за исключением группы, в которой меньше 5 элементов и ещё одной группы, содержащей сам элемент x. Таким образом получаем, что количество элементов больших x не менее $\frac{3N}{10}$. \Endproof

\Alg
\begin{enumerate}
    \item Найдем медиану медиан массива.
    \item Возьмем ее в качестве pivot.
    \item Проведем Partition.
    \item Если индекс опорного элемента i окажется больше данного k, то нам надо искать слева k-ю порядковую, в случае равенства завершаемся, в противном случае надо искать (i - k - 1)-ю порядковую.
\end{enumerate}

\Statement Данный алгоритм работает в худшем случае за $O(N)$ времени.

\Proof У нас есть три составляющих работы алгоритма на каждом шаге:
\begin{enumerate}
    \item Время на разделение массива на пятерки и сортировка каждой из них: $cN$.
    \item Время на поиск медианы медиан $T(\frac{N}{5})$.
    \item Время на поиск k-й порядковой не превзойдет времени его поиска в большей доле, то есть $T(\frac{7N}{10})$.
\end{enumerate}

Таким образом $T(N)\,=\,T(\frac{N}{5})+\,T(\frac{7N}{10})\,+\,cN$

По индукции: $T(N) \leq 10cN$

Следовательно $T(N)\,=\,T(\frac{N}{5})+\,T(\frac{7N}{10})\,+\,cN \leq \frac{10cN}{5} +\frac{7 \cdot 10cN}{10}\,+\,cN\,=\,10cN$ $\Rightarrow O(N)$
\Endproof

\pagebreak